
\section*{Общая характеристика работы}

\newcommand{\actuality}{\underline{\textbf{\actualityTXT}}}
\newcommand{\progress}{\underline{\textbf{\progressTXT}}}
\newcommand{\aim}{\underline{{\textbf\aimTXT}}}
\newcommand{\tasks}{\underline{\textbf{\tasksTXT}}}
\newcommand{\novelty}{\underline{\textbf{\noveltyTXT}}}
\newcommand{\influence}{\underline{\textbf{\influenceTXT}}}
\newcommand{\methods}{\underline{\textbf{\methodsTXT}}}
\newcommand{\defpositions}{\underline{\textbf{\defpositionsTXT}}}
\newcommand{\reliability}{\underline{\textbf{\reliabilityTXT}}}
\newcommand{\probation}{\underline{\textbf{\probationTXT}}}
\newcommand{\contribution}{\underline{\textbf{\contributionTXT}}}
\newcommand{\publications}{\underline{\textbf{\publicationsTXT}}}


{\actuality} Обзор, введение в тему, обозначение места данной работы в
мировых исследованиях и~т.\:п.

Традиционный подход (построение аналитических моделей):
\begin{itemize}
\item  сложная система разбивается на несколько независимых уровней
\item  определяются показатели эффективности каждого уровня
\item  производится построение, валидация и верификация <аналитических> моделей каждого уровня
\item  показатели каждого уровня оптимизируются в модельных сценариях
\item  производится подстройка модели
\end{itemize}


Предлагаемый подход (на основе машинного обучения, генерация фич):
\begin{itemize}
\item  сложная система данных воспринимается как "серый ящик"
\item  определяются показатели эффективности системы в целом 
\item  определяется конечное множество состояний системы
\item  из базовых признаков состояния системы производится порождение составных признаков 
\item  вырабатывается решающее правило
\item  решающее правило и модель принятия решения валидируется экспертами (трактовка результатов)
\end{itemize}

В данной работе предлагается рассмотреть несколько примеров использования машинного обучения в задачах управления инфраструктурой. Далее предлагается обобщить опыт и предложить метод построения и трактовки моделей, полученных при помощи алгоритмов машинного обучения в условиях современной инфраструктуры.


% разнородная по всем
% алгоритмов полно
% основные задачи 



% , можно использовать ссылки на другие
% работы~\cite{Gosele1999161} (если их нет, то в автореферате
% автоматически пропадёт раздел <<Список литературы>>). Внимание! Ссылки
% на другие работы в разделе общей характеристики работы можно
% использовать только при использовании \verb!biblatex! (из-за технических
% ограничений \verb!bibtex8!. Это связано с тем, что одна и та же
% характеристика используются и в тексте диссертации, и в
% автореферате. В последнем, согласно ГОСТ, должен присутствовать список
% работ автора по теме диссертации, а \verb!bibtex8! не умеет выводить в одном
% файле два списка литературы).

% {\progress} 
% Этот раздел должен быть отдельным структурным элементом по
% ГОСТ, но он, как правило, включается в описание актуальности
% темы. Нужен он отдельным структурынм элемементом или нет ---
% смотрите другие диссертации вашего совета, скорее всего не нужен.

{\aim} данной работы является \ldots

Для~достижения поставленной цели необходимо было решить следующие {\tasks}:
\begin{enumerate}
  \item Обзор и анализ существующих алгоритмов управления инфраструктурой на основе оперативных массивов данных
  \item Разработка и исследование методов и алгоритмов управления радио-ресурсами в телекоммуникационных сетях четвертого поколения (4G)
  \item Разработка и исследование методов и алгоритмов управления приоритетом обработки инцидентов на основе оперативных данных устройств железнодорожной автоматики и телемеханики
  \item Экспериментальное исследование характеристик разработанных моделей, обобщение на более широкий класс задач

\end{enumerate}

{\novelty}
\begin{enumerate}
  \item Впервые \ldots
  \item Впервые \ldots
  \item Было выполнено оригинальное исследование \ldots
\end{enumerate}

{\influence} \ldots

{\methods} Для решения задач, поставленных в работе, были использованы основные положения системного анализа, теории информации, теории вероятностей; для проектирования программной системы – методы объектно-ориентированного проектирования и язык UML; для программной реализации алгоритмов и системы – методы структурного, объектно-ориентированного и параллельного программирования.

{\defpositions}
\begin{enumerate}
  \item Способ формирования модели планирования ресурсов телекоммуникационной инфраструктуры, пригодной для применения методов машинного обучения. Децентрализованный метод управления радио-ресурсами в телекоммуникационных сетях четвертого поколения в целях понижения общего уровня интерференции и максимизации суммарной пропускной способности в сети базовых станций.
  \item Алгоритм и режимы работы элементов телекоммуникационной инфраструктуры сетей четвертого поколения, учитывающие требования по  качеству обслуживания пользователей.
  \item Алгоритмов и архитектура управления приоритетом обработки инцидентов на основе оперативных данных устройств железнодорожной автоматики и телемеханики.
  \item Алгоритмов и архитектура управления приоритетом обработки инцидентов на основе оперативных данных устройств железнодорожной автоматики и телемеханики.
\end{enumerate}
% В папке Documents можно ознакомиться в решением совета из Томского ГУ
% в файле \verb+Def_positions.pdf+, где обоснованно даются рекомендации
% по формулировкам защищаемых положений. 

{\reliability} полученных результатов обеспечивается \ldots \ Результаты находятся в соответствии с результатами, полученными другими авторами.


% Внедрение

Программная реализация системы  внедрена... \todo{ссылка на отчет о внедрении}
Программная реализация системы автоматического ранжирования инцидентов внедрена \todo{ссылка на отчет о результатах подконтрольной эксплуатации}

{\probation}
Основные результаты работы докладывались~на:
перечисление основных конференций, симпозиумов и~т.\:п.

{\contribution} Автор принимал активное участие \ldots

%\publications\ Основные результаты по теме диссертации изложены в ХХ печатных изданиях~\cite{Sokolov,Gaidaenko,Lermontov,Management},
%Х из которых изданы в журналах, рекомендованных ВАК~\cite{Sokolov,Gaidaenko}, 
%ХХ --- в тезисах докладов~\cite{Lermontov,Management}.

\ifnumequal{\value{bibliosel}}{0}{% Встроенная реализация с загрузкой файла через движок bibtex8
    \publications\ Основные результаты по теме диссертации изложены в XX печатных изданиях, 
    X из которых изданы в журналах, рекомендованных ВАК, 
    X "--- в тезисах докладов.%
}{% Реализация пакетом biblatex через движок biber
%Сделана отдельная секция, чтобы не отображались в списке цитированных материалов
    \begin{refsection}%
        \printbibliography[heading=countauthornotvak, env=countauthornotvak, keyword=biblioauthornotvak, section=1]%
        \printbibliography[heading=countauthorvak, env=countauthorvak, keyword=biblioauthorvak, section=1]%
        \printbibliography[heading=countauthorconf, env=countauthorconf, keyword=biblioauthorconf, section=1]%
        \printbibliography[heading=countauthor, env=countauthor, keyword=biblioauthor, section=1]%
        \publications\ Основные результаты по теме диссертации изложены в \arabic{citeauthor} печатных изданиях\nocite{qos,globecom}, 
        \arabic{citeauthorvak} из которых изданы в журналах, рекомендованных ВАК\nocite{globecom,ent}, 
        \arabic{citeauthorconf} "--- в тезисах докладов\nocite{globecom,ent}.
    \end{refsection}
}
При использовании пакета \verb!biblatex! для автоматического подсчёта
количества публикаций автора по теме диссертации, необходимо
их здесь перечислить с использованием команды \verb!\nocite!.
    

 % Характеристика работы по структуре во введении и в автореферате не отличается (ГОСТ Р 7.0.11, пункты 5.3.1 и 9.2.1), потому её загружаем из одного и того же внешнего файла, предварительно задав форму выделения некоторым параметрам

%Диссертационная работа была выполнена при поддержке грантов ...

%\underline{\textbf{Объем и структура работы.}} Диссертация состоит из~введения, четырех глав, заключения и~приложения. Полный объем диссертации \textbf{ХХХ}~страниц текста с~\textbf{ХХ}~рисунками и~5~таблицами. Список литературы содержит \textbf{ХХX}~наименование.

%\newpage
\section*{Содержание работы}
Во \underline{\textbf{введении}} обосновывается актуальность
исследований, проводимых в рамках данной диссертационной работы,
приводится обзор научной литературы по изучаемой проблеме,
формулируется цель, ставятся задачи работы, излагается научная новизна
и практическая значимость представляемой работы. В последующих главах
сначала описывается общий принцип, позволяющий ... , а потом идёт
апробация на частных примерах: ...  и ... .


\underline{\textbf{Первая глава}} посвящена ...

%  картинку можно добавить так:
% \begin{figure}[ht] 
%   \center
%   \includegraphics [scale=0.27] {latex}
%   \caption{Подпись к картинке.} 
%   \label{img:latex}
% \end{figure}

% Формулы в строку без номера добавляются так:
% \[ 
%   \lambda_{T_s} = K_x\frac{d{x}}{d{T_s}}, \qquad
%   \lambda_{q_s} = K_x\frac{d{x}}{d{q_s}},
% \]

\underline{\textbf{Вторая глава}} посвящена исследованию 

\underline{\textbf{Третья глава}} посвящена исследованию 

В \underline{\textbf{четвертой главе}} приведено описание 

В \underline{\textbf{заключении}} приведены основные результаты работы, которые заключаются в следующем:
%% Согласно ГОСТ Р 7.0.11-2011:
%% 5.3.3 В заключении диссертации излагают итоги выполненного исследования, рекомендации, перспективы дальнейшей разработки темы.
%% 9.2.3 В заключении автореферата диссертации излагают итоги данного исследования, рекомендации и перспективы дальнейшей разработки темы.
применение машинного обучения приносит свои плоды, однако внедрение затруднено тем, что клиенты не готовы доверять автоматизированным системам такого уровня абстракций. Трактовка работы моделей на базе машинного обучения и визуализация результата является необходимой частью успешного проекта. При этом желательно, чтобы трактовка происходила в реальном времени, с привлечением экспертов рабочей группы с самого начала разработки системы.
Опыт применения методов машинного обучения к задачам управления инфраструктурой российских железных дорог на Московской железной дороге  позволяет  ожидать  положительные  практические результаты при дальнейшем развитии и тиражировании достигнутых совместной рабочей группой специалистов.
По  результатам  доклада  генерального  директора  компании  ООО «Телум»  Павла  Бойко, Ученый  совет  ОАО «РЖД»  признал  перспективным применение современных методов машинного обучения к задачам управления инфраструктурой российских железных дорог и порекомендовал продолжить развитие, разработку и внедрение си-стемы в составе более широкой совместной груп-пы, с вовлечением отраслевых институтов и университетов (в т.ч. АО«ВНИИЖТ»). Представленные наработки в области интеллектуального анализа данных также рекомендовано апробировать в других областях железнодорожного транспорта.
\begin{enumerate}
  \item \todo{На основе анализа \ldots}
  \item \todo{Численные исследования показали, что \ldots}
  \item \todo{Математическое моделирование показало \ldots}
  \item \todo{Для выполнения поставленных задач был создан \ldots}
\end{enumerate}



%\newpage
% При использовании пакета \verb!biblatex! список публикаций автора по теме
% диссертации формируется в разделе <<\publications>>\ файла
% \verb!../common/characteristic.tex!  при помощи команды \verb!\nocite! 

\ifdefmacro{\microtypesetup}{\microtypesetup{protrusion=false}}{} % не рекомендуется применять пакет микротипографики к автоматически генерируемому списку литературы
\ifnumequal{\value{bibliosel}}{0}{% Встроенная реализация с загрузкой файла через движок bibtex8
  \renewcommand{\refname}{\large \authorbibtitle}
  \nocite{*}
  \insertbiblioauthor                          % Подключаем Bib-базы
  %\insertbiblioother   % !!! bibtex не умеет работать с несколькими библиографиями !!!
}{% Реализация пакетом biblatex через движок biber
  \insertbiblioauthor                          % Подключаем Bib-базы
  \insertbiblioother
}
\ifdefmacro{\microtypesetup}{\microtypesetup{protrusion=true}}{}

